% ch-a1-01.tex —— 第 1 篇（Dat）：PDF 18–27 / 印刷 5–14

\ysrc{5}因此，LLIF 是 CLLC 的一个推论；而由于从细模空间到粗模空间的过渡极具破坏性，LLIF 先验地只是对 CLLC 的一个远为粗略的近似。不过值得注意的是，Hellmann 的一个猜想提示了如何从 LLIF 再加上一点表示论来得到 CLLC，至少在 $G=\mathrm{GL}_n$、系数 $\Lambda=\overline\Q_\ell$ 的情形如此。事实上，注意 LLIF 给出了一个函子。确切地说，它在 $Z^1/\hatG$ 上给出拟凝聚层 $\wt W:=W\otimes_{\OO(Z^1)^{\hatG}}\OO(Z^1)$，带有 $G(F)$ 的作用，从而给出函子 $V\mapsto V\otimes^L_{G(F)}\wt W$。这当然不是正确的函子，但 Hellmann 的猜想 [Hel23, §4] 是：CLLC 应当具有形状 $V\mapsto V\otimes^L_{G(F)}\wt V$，其中 $\wt V$ 是 $\wt W$ 的某个商，由 Emerton 与 Helm 在他们关于 LLIF 的工作中定义。不过要注意，这一配方最朴素的推广对更一般的群并不够用，哪怕仅仅因为存在非通用（non generic）的尖点表示。

本文的目的是解释模空间 $Z^1/\hatG$ 的构造与主要性质，并特别强调粗模空间 $Z^1\sslash\hatG$。我们的第一项任务是弄清我们要使用哪一类 Langlands 参数。

\section{Langlands 参数。什么参数？}

记 ${}^LG:=\hatG\rtimes W_F$ 为 $G$ 的 Langlands 对偶。$L$-同态 $W_F\to{}^LG(\Lambda)$ 是第二投影的一个截面的群同态，故它具有形状 ${}^L\varphi(w)=(\varphi(w),w)$，其中 $\varphi:W_F\to\hatG(\C)$ 是某个 1-上闭链。${}^L\varphi$ 与 $\varphi$ 这二者互相唯一决定。

\subsection{复情形}

最初的经典局部 Langlands 对应预测一个有限到一的映射 $\pi\mapsto\varphi_\pi$，从 $\Rep_\C(G(F))$ 的不可约对象映到某个「Langlands 参数」集合。后者是文献中以两种不同风味出现的某些对象的 $\hatG(\C)$-共轭类：

\begin{enumerate}[label=(\arabic*),leftmargin=2.6em,itemsep=0.2em]
\item 映射 $\rho:W_F\times\mathrm{SL}_2(\C)\to{}^LG(\C)$，它在第一个因子上的限制是核为开的 $L$-同态，在第二个因子上的限制是代数的；或
\item 对 $(r,N)$，其中 $r:W_F\to{}^LG(\C)$ 是核为开的 $L$-同态，而 $N\in\mathrm{Lie}(\hatG_\C)$ 是幂幺元素，且满足 $\mathrm{Ad}_{r(w)}N=|w|N$。
\end{enumerate}

其实这里还漏掉一个重要的要求，文献中称为 Frobenius 半单性。它要求：对任意 $w\in W_F$，元素 $\rho(w,1)$（相应地 $r(w)$）在 ${}^LG(\C)$ 中是半单的。我们略去它的理由是：这一条件在族中表现欠佳，并且会妨碍我们定义一个参数化这些对象的良好簇（试想 $\hatG(\C)$ 中半单元素之集：它构成可构造子集，却不是局部闭子簇）。反之，当我们取粗模空间时，半单性会重新回来（试想 Chevalley–Steinberg 同构 $\hatG\sslash\hatG=\hat T\sslash W$）。

固定一个开紧子群 $K\subset W_F$。容易看出，对 $i=1,2$，在 $K$ 上平凡的 $(i)$ 型对象由某个仿射复簇 $Z^1_{(i)}$ 参数化。此外，有从 $(1)$ 型对象到 $(2)$ 型对象的映射 $\rho\mapsto(r,N)$，它由
\[
r(w)=\rho\left(w,\begin{pmatrix}|w|^{\frac12}&0\\0&|w|^{-\frac12}\end{pmatrix}\right),\qquad
N=d\rho\begin{pmatrix}0&1\\0&0\end{pmatrix}
\]
给出。该映射诱导簇的态射 $Z^1_{(1)}\to Z^1_{(2)}$。

\ysrc{6}由 Jacobson–Morozov 定理的一个著名应用，上述映射在 Frobenius 半单对象的 $\hatG(\C)$-共轭类上诱导双射，从而确保两种相应的 Langlands 参数概念一致。然而，该映射在一般的 $\hatG(\C)$-共轭类上不是单射，甚至不是满射（见 [MIY22, Example 3.5]）。相应地，$Z^1_{(1)}$ 与 $Z^1_{(2)}$ 的几何差异极大。

为了对究竟发生了什么有个概念，取 $G=\mathrm{GL}_2$、$K=I_F$。此时 $Z^1_{(1)}$ 是两个 4 维光滑分支的不交并 $\mathrm{GL}_2\sqcup(\C^\times\times\mathrm{PGL}_2)$，这两个分支分别对应 $\rho|_{\mathrm{SL}_2}$ 平凡与否。相反，$Z^1_{(2)}$ 连通，其两个 4 维不可约分支分别由 $N=0$ 的闭轨道（同构于 $\mathrm{GL}_2$）与 $N\ne0$ 的开轨道之闭包（该开轨道同构于 $\C^\times\times\mathrm{PGL}_2$）给出。这两个分支沿 3 维闭子簇相交，该闭子簇由形如 $\{(F,0)\}$、其中 $F$ 的两个特征值之比为 $q$ 的点组成。于是在此例中，$Z^1_{(1)}$ 表现为 $Z^1_{(2)}$ 的正规化的一个开子簇。更一般地，[MIY22] 证明映射 $Z^1_{(1)}\to Z^1_{(2)}$ 分离不可约分支，并且是「弱双有理的」（weakly birationnal\tn{底本此处作「birationnal」（双 n），疑为 birational 之误。}）。

考察相应粗模空间上诱导的映射 $Z^1_{(1)}\sslash\hatG\to Z^1_{(2)}\sslash\hatG$ 是有意思的。作为一般原则，$Z^1_{(i)}\sslash\hatG$ 的复点对应 $Z^1_{(i)}(\C)$ 的闭轨道（见第 5 节），而这些闭轨道是 $Z^1_{(i)}(\C)$ 中「半稳定」对象的轨道。应用 Hilbert–Mumford 判别法可知，$Z^1_{(1)}(\C)$ 的半稳定点恰为 Frobenius 半单对象，而 $Z^1_{(2)}(\C)$ 的半稳定点是形如 $(r,0)$、其中 $r$ 为 Frobenius 半单的点。换言之，$(Z^1_{(1)}\sslash\hatG)(\C)$ 恰为 Langlands 参数之集，而 $(Z^1_{(2)}\sslash\hatG)(\C)$ 是所谓「无穷小特征标」之集（Vogan 的术语），并且映射 $Z^1_{(1)}\sslash\hatG\to Z^1_{(2)}\sslash\hatG$ 把 Langlands 参数送到它的无穷小特征标。

这段讨论的要点是：在系数 $\Lambda=\C$ 的情形，对于设想中的 CLLC 所用的模空间 $Z^1/\hatG$，先验地有两个自然候选。尽管 $Z^1_{(1)}$ 看起来更接近 Langlands 参数的经典概念，它却不是 CLLC 的合理候选。事实上，所期望的函子必须把平凡表示与 Steinberg 表示送到某些不可分解层。既然不可分解，这些层就应当支撑在单个分支上。由与经典 LLC 的相容性，联系于 Steinberg 的层应落在 $\rho|_{\mathrm{SL}_2}$ 非平凡的那个分支上，而联系于平凡表示的层落在另一分支上。但这两个层没有非平凡扩张，而这两个表示却有，因此这与所期望函子的全忠实性不相容。事实上，[Hel23] 与 [BZCHN20] 中出现的模空间都是 $Z^1_{(2)}/\hatG$。相应地，与 LLIF 相关的粗模空间是 $Z^1_{(2)}\sslash\hatG$。

\subsection{$\ell$-进情形}

把上述讨论从 $\C$ 搬到 $\overline\Q_\ell$：或者重复定义，或者选取域同构 $\C\simeq\overline\Q_\ell$。[Hel20b] 中原来的 LLIF 问题在系数 $\Lambda$ 取 $\Z_\ell$ 的有限扩张时提出，[Zhu25] 与 [FS21] 中设想的 CLLC 亦然。这意味着我们需要 $\Z_\ell$ 上的模空间 $Z^1/\hatG$。

\ysrc{7}注意上面 $(2)$ 型对象 $(r,N)$ 在把 $\C$ 换成任意 $\Z_\ell$-代数 $R$ 后完全有意义，并由 $\Z_\ell$ 上的仿射概形参数化，从而给出 $Z^1_{(2)}$ 的一个整模型。然而 LLIF 的一项要求是：几何特殊纤维 $Z^1(\overline\F_\ell)$ 应当参数化连续伽罗瓦表示（即 $L$-同态）$\overline\rho:\Gamma_F\to{}^LG(\overline\F_\ell)$。于是需要一套从对 $(r,N)$ 得到这样的 $\overline\rho$ 的配方。作者所知唯一的配方要用到从 $\mathrm{Lie}\,\hatG_{\F_\ell}$ 的幂幺元素到 $\hatG(\overline\F_\ell)$ 的幂幺元素的指数映射，所以第一个问题是：当 $\ell$ 较小时（例如对 $\mathrm{GL}_n$ 有 $\ell<n$ 时），这样的指数映射可能不存在。无论如何，假设 $\ell$ 大到足以有指数映射，配方如下：固定一个平凡化 $\Z_\ell(1):=\varprojlim_n\mu_{\ell^n}\isom\Z_\ell$，并固定 $W_F$ 中 Frobenius 的一个提升 $\mathrm{Fr}$。于是得到态射 $t_\ell:I_F\to\Z_\ell$ 与收缩 $W_F\to I_F$，$w\mapsto i_w$。此时映射 $w\mapsto\overline\rho(w):=r(w)\exp(t_\ell(i_w)N)$ 定义了所要的 $L$-同态 $\Gamma_F\to{}^LG(\overline\F_\ell)$，并在 $\hatG(\overline\F_\ell)$-共轭类之间建立双射。同一配方实际上对任意完备局部 $\Z_\ell$-代数都有效，并允许比较 $(r,N)$ 与 $\overline\rho$ 的形变空间。正是在这里我们会遇到更多麻烦，如下例所示。

\begin{exam}{1.1}
设 $G=\mathrm{GL}_2$，$\ell>2$ 且 $q\equiv-1\pmod\ell$。取 $\zeta\in\mu_\ell(\overline\Q_\ell)$，并定义特征标 $\chi:I_F.\mathrm{Fr}^{2\Z}\to\overline\Z_\ell^\times$ 为 $\chi(i)=\zeta^{t_\ell(i)}$、$\chi(\mathrm{Fr}^2)=1$。则 $W_F$ 的 2 维不可约 $\overline\Q_\ell$-表示 $\rho:=\mathrm{ind}_{I_F.\mathrm{Fr}^{2\Z}}^{W_F}\chi$ 不可约，其半单化归约同构于平凡特征标与 $\overline\F_\ell$ 上二次非分歧特征标 $\varepsilon$ 的直和 $r=1\oplus\varepsilon$。由于 $\rho$ 不可约，它含有一个不变的 $\Z_\ell$-格 $L_\rho$，其归约 $\overline\rho$ 是 $\varepsilon$ 被 $1$ 的一个非平凡扩张。于是对应于 $\overline\rho$ 的 Weil–Deligne 表示具有形状 $(r,N)$，其中 $N\ne0$。但 $(r,N)$ 的任何提升都必须有非平凡的 $N$，可见 $(r,N)$ 的任何提升都不匹配 $L_\rho$；也就是说，$(r,N)$ 与 $\overline\rho$ 的形变空间不同。
\end{exam}

于是我们看到，$(2)$ 型对象一般不能给出 $\Z_\ell$ 上正确的模空间。不过上述配方在系数 $\overline\Q_\ell$ 上对任意 $\ell$ 都有效，并把上面的 $(2)$ 型对象送到 $(3)$ 型对象：

\begin{enumerate}[label=(\arabic*),start=3,leftmargin=2.6em]
\item $L$-同态 ${}^L\varphi:W_F\to{}^LG(\overline\Q_\ell)$，且为 $\ell$-进连续的。
\end{enumerate}

这里「$\ell$-进连续」是指：当给 $\hatG(\overline\Q_\ell)$ 配备由 $\overline\Q_\ell$ 上自然拓扑诱导的拓扑时，${}^L\varphi$ 连续。此外，Grothendieck 的拟幂幺性定理表明，这在相应的 $\hatG(\overline\Q_\ell)$-共轭类之间诱导双射。

类似地，若 $R$ 是 $\Z_\ell$ 的有限扩张，或者更一般地，是一个剩余域包含于 $\overline\F_\ell$ 的完备局部代数，我们也可以对 $L$-同态 ${}^L\varphi:W_F\to{}^LG(R)$ 施加 $\ell$-进连续性。在 $R=\overline\F_\ell$ 的情形，连续的 ${}^L\varphi$ 会唯一地延拓到 $\Gamma_F$。因此，如果我们能把「$\ell$-进连续」这一概念推广到任意 $\Z_\ell$-代数 $R$，并证明相应的函子
\[
R\mapsto\bigl\{{}^L\varphi:W_F\to{}^LG(R)\ \bigm|\ \text{${}^L\varphi$ 为 $\ell$-进连续的 $L$-同态}\bigr\}
\]
可由仿射 $\Z_\ell$-概形 $Z^1$ 表示，那么 $Z^1$ 的 $\overline\F_\ell$-点当然就参数化伽罗瓦表示（$L$-参数），而它们的完备局部环就是通常的加框架形变环，这正是 LLIF 纲领所需。

\ysrc{8}

\section{如何处理连续性}

当 $R$ 是任意 $\Z_\ell$-代数时，$L$-同态 ${}^L\varphi:W_F\to{}^LG(R)$ 为「$\ell$-进连续」有若干种可能的定义。

\subsection{Helm 的原始进路}

注意，若 $R$ 是 $\ell$-进分离的，则 $\hatG(R)$ 上有一个自然的分离群拓扑，它由归约映射 $\hatG(R)\to\hatG(R/\ell^mR)$ 的核生成。于是具体地说，在此情形称 ${}^L\varphi$ 为 $\ell$-进连续的，是指对每个 $m$，复合映射 $W_F\to\hatG(R/\ell^m)\rtimes W_F$ 在 $\hatG(R/\ell^m)$ 取离散拓扑、$W_F$ 取自然拓扑时连续。仿照 [Hel20b] 与 [DHKM20]，对任意的 $\Z_\ell$-代数，我们于是可以宣布 ${}^L\varphi$ 为 $\ell$-进连续的，如果它由某个分离 $\Z_\ell$-代数 $R_0$ 上的 $\ell$-进连续 $L$-同态经前推 $R_0\to R$ 得到。

这种定义的明显缺点是：完全不清楚函子 $R\mapsto\{\ell\text{-进连续的 }{}^{L}\varphi\}$ 为何应当可表示，甚至不清楚它是否定义 fppf 层。在眼下的情形，由于 $W_F$ 的结构\tn{底本此处作「struture」，疑为 structure 之误。}相当简单，结果会颇为神奇地确实如此。

\subsection{凝聚数学}

[FS21] 采取了另一条进路。那里把 $R$ 看作凝聚 $\Z_\ell$-代数 $S\mapsto R(S)$：对任意投射有限集 $S$，$R(S)$ 是映射 $S\to R$ 构成的环，这些映射要能经过某个有限生成 $\Z_\ell$-子模分解，并且关于该子模上的 $\ell$-进拓扑连续。相应地，$\hatG(R)$ 成为凝聚群。另一方面，拓扑群 $W_F$ 也可看作凝聚群，而 [FS21] 的想法是考虑凝聚群的 $L$-同态。事实上，这正是他们关于 Fargues–Fontaine 曲线上 Hecke 算子的上同调框架中自然出现的概念。

这个定义也可以用通常的拓扑语言表达。给 $R$ 配备最细的拓扑，使其在所有有限生成 $\Z_\ell$-子模上限制为自然的 $\ell$-进拓扑（在 [Zhu25] 中，这个拓扑称为「ind-$\ell$-进」）。这在任意有限秩自由 $R$-模上诱导 $\Z_\ell$-线性拓扑。于是宣布 ${}^L\varphi$ 连续，如果对任意（等价地，对某一个忠实）代数表示 $\hatG\to\mathrm{GL}(V)$，$\ker(W_F\to\Aut(\hatG))$ 在 $V\otimes R$ 上的诱导作用是连续的。这里的连续性指：任意 $v\in V\otimes R$ 都含于某个有限生成 $\Z_\ell$-模，该模在 $W_F$ 的某个开子群下稳定，且诱导作用是连续的。

先验地不清楚这个概念是否涵盖前一个。也就是说，$\ell$-进分离代数 $R$ 的 ind-$\ell$-进拓扑一般比 $\ell$-进拓扑更细。然而，同样由于 $W_F$ 的特殊结构，两个概念会给出同一个模空间。主要原因是 Frobenius 的魔力以及它与惯性的相互作用。

\subsection{驯顺惯性的离散化}

由于 $\ell\ne p$，无论对 ${}^L\varphi$ 取何种「$\ell$-进连续」定义，$\varphi$ 在野惯性子群 $P_F$ 上的限制都会经过一个有限商分解。因此这里唯一的拓扑困难来自驯顺惯性商 $I_F/P_F$。它是 pro-循环群，故取一个生成元 $s$。再取 $W_F/P_F$ 中 Frobenius 的一个提升 $\mathrm{Fr}$。它们满足关系 $\mathrm{Fr}\,s\,\mathrm{Fr}^{-1}=s^q$。于是考虑 $W_F/P_F$ 的子群
\[
\langle\mathrm{Fr},s\rangle=s^{\Z[\frac1q]}\rtimes\mathrm{Fr}^{\Z},
\]
记 $W^0_F$ 为它在 $W_F$ 中的原像，并给它配备这样的拓扑：它延拓 $P_F$ 的投射有限拓扑，并在 $\langle\mathrm{Fr},s\rangle$ 上诱导离散拓扑。

\ysrc{9}于是 $W^0_F$ 是 $\Gamma_F$ 的稠密子群，而 $\Gamma_F$ 同等于 $W^0_F$ 的投射有限完备化，因为 $W^0_F$ 的每个有限商都是 $F$ 的某个有限扩张的伽罗瓦群。注意（与 $W_F$ 作为 $\Gamma_F$ 的子群相对照\tn{底本此处作「the subgroup $W_F$ of $G_F$」，其中 $G_F$ 疑为 $\Gamma_F$ 之误。}），$W^0_F$ 作为 $W_F$ 的子群在很大程度上依赖于 $\mathrm{Fr}$ 与 $s$ 的选取。

现在宣布 $L$-同态 ${}^L\varphi:W^0_F\to{}^LG(R)$ 连续，如果它关于 ${}^LG(R)$ 上的离散拓扑与 $W^0_F$ 上的上述拓扑连续。这等价于要求相应的 1-上闭链 $\varphi:W^0_F\to\hatG(R)$ 在 $P_F$ 的某个开子群上平凡。

\subsection{一个候选的模空间}

固定一个在 $\hatG$ 上平凡作用的开正规子群 $P\subset P_F$，则容易看出 $\Z[\frac1p]$-代数上的函子
\[
R\mapsto\{\text{1-上闭链 }\varphi:W^0_F/P\to\hatG(R)\}
\]
可由有限生成 $\Z[\frac1p]$-代数 $R_{\hatG,1P}$ 表示。于是置
\[
Z^1=Z^1(W^0_F/P,\hatG):=\Spec(R_{\hatG,1P}).
\]
具体地说，离散群 $W^0_F/P$ 有限生成，设由 $n$ 个元素生成；而 $Z^1$ 是 $(\hatG)^n$ 中由这些生成元之间的关系与上闭链条件定义的闭子概形。

尽管 $Z^1$ 先验地依赖于 $W_F$ 的离散化 $W^0_F$ 之选取，我们仍有几项检验可做，以判断它是否是 LLIF 或 CLLC 的合理候选。

\begin{enumerate}[label=(\arabic*),leftmargin=2.6em,itemsep=0.3em]
\item 当 $R=\overline\F_\ell$，或 $R$ 是任意以 $\overline\F_\ell$ 为剩余域的诺特完备局部代数时，连续的 ${}^L\varphi:W^0_F\to{}^LG(R)$ 可唯一地连续延拓到 $W^0_F$ 的投射有限完备化，即 $\Gamma_F$。于是 $Z^1$ 满足要求：其 $\overline\F_\ell$-点参数化取值于 ${}^LG(\overline\F_\ell)$ 的伽罗瓦表示，且这些表示的形式邻域与它们通常的加框架形变空间一致。
\item 当 $R=\overline\Q_\ell$ 时，我们不能以同样的方式论证，但我们断言：任意连续 ${}^L\varphi:W^0_F\to{}^LG(\overline\Q_\ell)$ 都唯一地连续延拓到 $W_F$，从而 $Z^1$ 的 $\overline\Q_\ell$-点如所愿参数化上面 $(3)$ 型对象。注意，利用如 [Iwa55] 中的连续分裂 $W_F\simeq P_F\rtimes(W_F/P_F)$，只需在 $P=P_F$ 时验证此事。选取 ${}^LG$ 的一个忠实表示后，我们\tn{底本此处作「wa may assume」，疑为 we may assume 之误。}可以假设 $G=\mathrm{GL}_n$，于是 ${}^L\varphi$ 由态射 $\varphi:W^0_F/P_F=\langle\mathrm{Fr},s\rangle\to\mathrm{GL}_n(\overline\Q_\ell)$ 给出。此时关系 $\mathrm{Fr}\,s\,\mathrm{Fr}^{-1}=s^q$ 表明 $\varphi(s)$ 的特征值都是单位根，从而半单部分 $\varphi(s)_{\mathrm{ss}}$ 的阶与 $p$ 互素。另一方面，把幂幺部分写作 $\varphi(s)_{\mathrm{u}}=\exp(N)$，我们看到 $\varphi$ 通过 $I_F/P_F$ 上由
\[
\varphi(s^i)=(\varphi(s)_{\mathrm{ss}})^i\exp(t_\ell(i)N),\qquad i\in\Z[\tfrac1q],
\]
给出的公式延拓到 $W_F/P_F$。
\end{enumerate}

对更一般的 $\Z_\ell$-代数 $R$，下述定理蕴含任意连续 $L$-同态 $W^0_F\to{}^LG(R)$ 实际上都能延拓为唯一的 $\ell$-进连续 $L$-同态 $W_F\to{}^LG(R)$。

\ysrc{10}

\begin{thm}{2.1}
固定如上开子群 $P\subset P_F$，并记 ${}^L\varphi_{\hatG,1P}$ 为泛 $L$-同态 $W^0_F\to{}^LG(R_{\hatG,1P})$。
\begin{enumerate}[label=(\arabic*),leftmargin=2.6em,itemsep=0.2em]
\item 环 $R_{\hatG,1P}$ 是 $\ell$-进分离的。
\item 泛 ${}^L\varphi_{\hatG,1P}$ 唯一地延拓为 $\ell$-进连续 $L$-同态 ${}^L\varphi_{\hatG,1P,\ell}:W_F\to{}^LG(R_{\hatG,1P}\otimes\Z_\ell)$，且它对 $W_F$ 上在 $P$ 上平凡的 $\ell$-进连续 $L$-同态是泛的。
\item $L$-同态 ${}^L\varphi_{\hatG,1P,\ell}$ 也是 ind-$\ell$-进连续的（见 2.2），且它对 $W_F$ 上在 $P$ 上平凡的 ind-$\ell$-进连续 $L$-同态是泛的。
\end{enumerate}
\end{thm}

$(2)$ 的一个推论是：基变换 $Z^1(W^0_F/P,\hatG)_{\Z_\ell}$ 不依赖\footnote{目前尚不知道完整的概形 $Z^1(W^0_F/P,\hatG)$ 或叠 $Z^1(W^0_F/P,\hatG)/\hatG$ 是否真的依赖于离散化的选取，但至少由 [DHKM20, Thm. 4.18] 可知，粗商概形 $Z^1(W^0_F/P,\hatG)\sslash\hatG$ 与这一选取无关。}于我们对离散化 $W^0_F$ 的选取。$(3)$ 的要点是：这个基变换就是 [FS21, Chapter VIII] 中所考虑的凝聚 $L$-同态的模空间。我们将在接下来几节解释该定理的证明。

\subsection{关于导出结构的注}

上面定义的仿射概形 $Z^1$ 完全适用于 LLIF，但它是否足以应付 CLLC，先验地并不清楚。这是因为 CLLC 的凝聚一侧，即 $Z^1$ 上拟凝聚模的（导出 $\infty$）范畴，对 $Z^1$ 上可能的非平凡导出结构相当敏感。这样的非平凡导出结构可能出现，例如：在选取 $W^0_F/P$ 的一组 $n$ 个生成元后得到的闭浸入 $Z^1\hookrightarrow(\hatG)^n$ 若不正则，就会出现这种情况。技术上，应当在 $\Z[\frac1p]$-代数的某个「导出范畴」中工作，即「animated」$\Z[\frac1p]$-代数的 $\infty$-范畴，并像以前一样考虑同一个函子
\[
R\mapsto\{\text{1-上闭链 }\varphi:W^0_F/P\to\hatG(R)\},
\]
但配以正确的 $\infty$-范畴意义下的 1-上闭链概念。更多细节，以及「该函子确实可由仿射导出概形 $\wt Z^1$ 表示、其经典底概形为 $Z^1$」的证明，参见 [Zhu25, §2.2]。

尽管在研究整体伽罗瓦群的表示模空间时这一导出框架可能是必要的，但由 [Zhu25, Prop. 2.2.13] 与下述定理的第一条可知，在我们的情形有 $\wt Z^1=Z^1$。

\begin{thm}{2.2}
仿射概形 $Z^1(W_F/P,\hatG)$ 是 $\Z[\frac1p]$ 上维数为 $\dim\hatG$ 的约化平坦局部完全交。
\end{thm}

上述两个定理都蕴含 $Z^1(W_F/P,\hatG)$ 的任一不可约分支在 $\Z[\frac1p]$ 上忠实平坦。这两个定理的证明将在接下来两节给出，同时给出 $Z^1(W_F/P,\hatG)_{\Z[\frac1p]}$ 连通分支的参数化，以及用较小的群对这些分支的描述。

\section{玩具模型：驯顺参数}

本节聚焦 $P=P_F$ 且群 $G$ 驯顺分歧的情形，此时 $P_F$ 在 $\hatG$ 上的作用平凡。

实际上，我们将考虑一个稍有不同的设定。记 $\Lambda$ 为 $\Z[\frac1p]$ 的有限平坦整扩张，$H$ 为 $\Lambda$ 上的群概形。然后考虑
\ysrc{11}仿射 $\Lambda$-概形
\[
X_{H,1}:=\underline{\Hom}(W^0_t,H),
\]
它参数化从 $W^0_F$ 的驯顺商 $W^0_t:=W^0_F/P_F=\langle\mathrm{Fr},s\rangle$ 到 $H$ 的态射。具体地说，它是 $H\times H$ 的闭子概形，表示 $\Lambda$-代数上的函子
\[
R\mapsto\Hom(W^0_t,H(R))=\{(F,\sigma)\in H(R)\times H(R),\ F\sigma F^{-1}=\sigma^q\}.
\]

我们最关心的是 $H=\hatG\rtimes W$ 的情形，其中 $W$ 是 $W^0_t$ 的有限商，$W_F$ 在 $\hatG$ 上的作用经由它分解。此时我们最终关心的概形 $Z^1(W^0_t,\hatG)$ 是上述概形 $X_{H,1}$ 的直和项。不过，考虑更广的一类群概形 $H$ 会有用。

\subsection{广义约化群概形}

回忆底概形 $S$ 上的群概形 $H$ 称为约化的，如果它在 $S$ 上光滑，且所有几何纤维都是约化且连通的。允许一定非连通性的一点推广会很有用。为此回忆：对 $S$ 上任意光滑群概形，存在唯一的开子群概形 $H^0\subset H$，其几何纤维是 $H$ 的几何纤维的单位分支。此外，在 $H^0$ 约化的假设下，已知 [Con14, Prop. 3.1.3]：商层 $\pi_0(H):=H/H^0$ 可由 $S$ 上分离 étale 群概形表示。

称 $H$ 为广义约化的（generalized reductive），如果它满足以下性质。
\begin{itemize}[leftmargin=1.6em,itemsep=0.15em]
\item $H$ 是 $S$ 上的光滑仿射群概形；
\item $H^0$ 是 $S$ 上的约化群概形；
\item $\pi_0(H):=H/H^0$ 是 $S$ 上的有限（étale）群概形。
\end{itemize}
此外，称这样的 $H$ 是分裂的，如果 $H^0$ 是分裂约化群概形且 $\pi_0(H)$ 常值。我们需要的关于这些对象的主要性质概括在下面两个引理中。

\begin{lem}{3.1}
设 $H$ 在 $S$ 上分裂，$(B,T,X)$ 是 $H^0$ 的一个固定数据（pinning）。
\begin{enumerate}[label=(\arabic*),leftmargin=2.6em,itemsep=0.2em]
\item 环面 $T$ 的正规化子 $N_H(T)$、相应地 Borel 对 $(B,T)$ 的正规化子 $N_H(T,B)$，都是分裂广义约化的，单位分支为 $T$。
\item 若 $H^0$ 的中心 $Z(H^0)$ 在 $S$ 上光滑，则固定数据的正规化子 $N_H(B,T,X)$ 是广义约化的，单位分支为 $Z(H^0)^0$。
\end{enumerate}
\end{lem}

注意，对上述引理中每个群概形 $N$，都在几何纤维上验证：自然态射 $N\to\pi_0(H)$ 是满的（即在 étale 层意义下满），而核 $N\cap H^0$ 属于约化群概形的理论，其有效参考文献是 [Con14]。

\begin{lem}{3.2}
设 $S=\Spec R$，其中 $R$ 是数域的正规子环。
\begin{enumerate}[label=(\arabic*),leftmargin=2.6em,itemsep=0.2em]
\item $S$ 上任意广义约化 $H$ 都在某个有限整扩张上分裂。
\item 若 $H$ 在 $S$ 上广义约化，且 $P\subset H(S)$ 是可解有限子群、其阶在 $S$ 上可逆，则 $P$ 的中心化子 $C_H(P)$ 与正规化子 $N_H(P)$ 都是广义约化的，且 $C_H(P)^0=N_H(P)^0$。
\end{enumerate}
\end{lem}

$(1)$ 的证明见 [DHKM20, Theorem 1.13]，$(2)$ 的证明见 [DHKM20, Theorem A.12]。

\ysrc{12}

\subsection{平坦性与完全交}

以下我们总假设 $\Lambda$ 上的群概形 $H$ 是「广义约化」的。

\begin{prop}{3.3}
$X_{H,1}$ 在 $S:=\Spec\Lambda$ 上是合冲的（syntomic，即平坦且局部完全交），相对维数与 $H$ 在 $S$ 上的相对维数相同。
\end{prop}

\begin{proof}
有笛卡尔方块
\begin{center}
\begin{tikzcd}
X_{H,1} \arrow[r] \arrow[d] & H\times_S H \arrow[d] \\
S \arrow[r] & H
\end{tikzcd}
\end{center}
其中右侧竖直映射为 $H\times_S H\to H$，$(F,\sigma)\mapsto F\sigma F^{-1}\sigma^{-q}$。这表明光滑 $S$-概形 $H\times_S H$ 的闭子概形 $X_{H,1}$ 由 $\dim_S H$ 个函数的零点定义。因此，为证明映射 $X_{H,1}\to S$ 是合冲的（甚至是整体完全交），只需证明其所有几何纤维的维数等于 $\dim_S(H\times_S H)-\dim_S H=\dim_S H$（其中 $\dim_S$ 表示在 $S$ 上的相对维数）。

于是设 $L$ 是 $\Lambda$ 上的代数闭域，为记号简便记 $X:=X_{H,1}$。由于 $X_L$ 是态射 $H_L\times H_L\to H_L$、$(F,\sigma)\mapsto F\sigma F^{-1}\sigma^{-q}$ 在 $1$ 处的纤维，$X_L$ 的每个不可约分支维数 $\ge\dim H_L=\dim_S H$。另一方面，设 $\sigma\in H(L)$ 在第二投影 $\pi_2:X\to H$ 的像中。则纤维 $\pi_2^{-1}(\sigma)\subset X$ 是 $H_L$ 中 $\sigma$ 的中心化子 $H_\sigma$ 上的一个主齐性空间（torsor），故维数为 $\dim H_\sigma$。由此，$\sigma$ 的 $H^0_L$-共轭类 $\tilde\sigma$ 的原像 $\pi_2^{-1}(\tilde\sigma)$ 的维数为 $\dim H_\sigma+(\dim H_L-\dim H_\sigma)=\dim H_L=\dim_S H$。于是所需的等维性可由下述断言推出：$\pi_2(X(L))$ 由 $H(L)$ 中有限多个共轭类组成。

为证明该断言，设 $\sigma\in\pi_2(X(L))$，考虑其 Jordan 分解 $\sigma=\sigma_{\mathrm{ss}}\sigma_{\mathrm{u}}$。我们将证明 $\sigma_{\mathrm{ss}}$ 的阶有与 $\sigma$（以及 $L$）无关的有限上界。由于有界有限阶半单元素的共轭类只有有限多个，而此类元素的中心化子中幂幺元素的共轭类也只有有限多个 [FG12]，断言得证。于是设 $e$ 为 $\pi_0(H_L)$ 的阶。则 $(\sigma_{\mathrm{ss}})^e\in H^0_L$ 属于 $H_L$ 的某个极大环面 $T$。由于 $(\sigma_{\mathrm{ss}})^{qe}$ 在 $H(L)$ 中共轭于 $(\sigma_{\mathrm{ss}})^e$，且也属于 $T$，存在 $n\in N_{H(L)}(T)$ 使得 $n(\sigma_{\mathrm{ss}})^e n^{-1}=(\sigma_{\mathrm{ss}})^{qe}$。记 $\omega$ 为 $H^0_L$ 的 Weyl 群的阶，则得 $(\sigma_{\mathrm{ss}})^e=n^{e\omega}(\sigma_{\mathrm{ss}})^e n^{-e\omega}=(\sigma_{\mathrm{ss}})^{q^{e\omega}e}$，从而 $(\sigma_{\mathrm{ss}})^{e(q^{e\omega}-1)}=1$。
\end{proof}

\begin{rem}{3.4}
由证明可知：$(X_{H,1})_L$ 的任一不可约分支都是 $\pi_2(X_{H,1}(L))$ 中某个 $H^0$-共轭类的原像之闭包。换言之，任一不可约分支都是映射
\[
H^0_L\times(H_\sigma)^0\longrightarrow(X_{H,1})_L,\qquad (h,k)\mapsto(hFkh^{-1},h\sigma h^{-1})
\]
的像之闭包，其中 $(F,\sigma)\in X(L)$ 取定。
\end{rem}

\begin{rem}{3.5}
$H^0$ 是约化的这一事实，在利用幂幺类集合的有限性时至关重要。没有这一假设，模空间可能不是局部完全交（l.c.i.），从而其导出类比（见第 2.5 小节）可能有非平凡的高阶结构。例如当 $H$ 是 Borel 子群时就会出现这种情况，见 [Hel23, Example 2.3]。
\end{rem}

\ysrc{13}

\subsection{无阻碍点与一般光滑性}

设 $L$ 是 $\Lambda$ 上的代数闭域，$x\in X_{H,1}(L)$ 是 $X_{H,1}$ 的一个 $L$-点。我们计算切空间 $T_x(X_{H,1})_L$ 及其作为 $L$-向量空间的维数。为略微简化记号，本节记 $X$ 为 $X_{H,1}$。

\begin{lem}{3.6}
设 $\varphi:W^0_t\to H(L)$ 是联系于 $x$ 的态射，并设 $W^0_t$ 经由 $\varphi$ 与伴随作用的复合作用在 Lie 代数 $\mathrm{Lie}\,H$ 上。
\begin{enumerate}[label=(\arabic*),leftmargin=2.6em,itemsep=0.2em]
\item $T_xX_L\simeq Z^1(W^0_t,\mathrm{Lie}\,H_L)$。
\item $\dim_L T_xX_L=\dim H_L+\dim_L H^0\bigl(W^0_t,(\mathrm{Lie}\,H_L)^*\otimes\omega\bigr)$，
\end{enumerate}
其中 $\omega$ 是把 $\mathrm{Fr}$ 送到 $q$、把 $s$ 送到 $1$ 的「分圆特征标」。特别地，$x$ 是 $X_L$ 的光滑点，当且仅当 $H^0\bigl(W^0_t,(\mathrm{Lie}\,H_L)^*\otimes\omega\bigr)=0$。
\end{lem}

\begin{proof}
$(1)$ 方便的做法是把 $X$ 看作 $H^{W^0_t}$ 中由条件 $\varphi(ww')\varphi(w)^{-1}\varphi(w')^{-1}=1$ 定义的闭子概形。此时 $T_xX_L$ 同等于相应线性映射 $\prod_w T_{\varphi(w)}H_L\to T_1H_L=\mathrm{Lie}\,H_L$ 的公共核。经右平移把 $T_hH_L$ 与 $\mathrm{Lie}\,H_L$ 等同，即把 $T_xX_L$ 实现为若干线性映射 $\prod_w\mathrm{Lie}\,H_L\to\mathrm{Lie}\,H_L$ 的公共核。为计算这些映射，回忆切空间 $T_xX_L$ 可看作映射 $X(L[\varepsilon])\xrightarrow{\pi}X(L)$ 在 $x$ 处的纤维。于是 $T_xX_L$ 的元素 $\tilde x$ 由某个同态 $\tilde\varphi:W^0_t\to H(L[\varepsilon])$ 给出，使得对任意 $w$ 有 $\tilde\varphi(w)\in T_{\varphi(w)}H_L$。在上述右平移下，这对应于 $\prod_w\mathrm{Lie}\,H_L$ 中的元素 $(\tilde h_w:=\tilde\varphi(w)\varphi(w)^{-1})_{w\in W}$。条件 $\tilde\varphi(ww')\tilde\varphi(w)^{-1}\tilde\varphi(w')^{-1}=1$ 于是化为等式\tn{底本此处「The condition $\tilde\varphi(ww')\tilde\varphi(w)^{-1}\tilde\varphi(w')^{-1}=1$ then」一句重复排版（断行处的重复），译文只保留一次。}
\[
\tilde h_{ww'}=\tilde\varphi(ww')\varphi(ww')^{-1}=\tilde\varphi(w)\tilde\varphi(w')\varphi(w')^{-1}\varphi(w)^{-1}=h_w\varphi(w)h_{w'}\varphi(w)^{-1}.
\]
用加性记号，这意味着 $\tilde h_{ww'}=h_w+\mathrm{Ad}_{\varphi(w)}(h_{w'})$，这正是 $(1)$ 所述的上闭链条件。

$(2)$ 可以用 [DHKM20, Lemma 5.1] 中的显式计算证明。这里给出另一个论证。有正合列
\[
0\to(\mathrm{Lie}\,H_L)^{\varphi(W^0_t)}\to\mathrm{Lie}\,H_L\xrightarrow{\ \partial\ }Z^1(W^0_t,\mathrm{Lie}\,H_L)\to H^1(W^0_t,\mathrm{Lie}\,H_L)\to0,
\]
其中 $\partial\tilde h=(w\mapsto\mathrm{Ad}_{\varphi(w)}\tilde h-\tilde h)$。这表明
\[
\dim_L T_xX_L=\dim H_L+\dim_L H^1(W^0_t,\mathrm{Lie}\,H_L)-\dim_L H^0(W^0_t,\mathrm{Lie}\,H_L).
\]
现在，dévissage $s^{\Z[\frac1q]}\hookrightarrow W^0_t\twoheadrightarrow\mathrm{Fr}^{\Z}$ 让我们分阶段计算 $W^0_t$ 的上同调，即经由一个谱序列。由于 $\Z$ 与 $\Z[\frac1q]$ 在有限维 $L$-向量空间上的上同调维数为 1，且满足 Euler–Poincaré 维数的消没，我们看到 $W^0_t$ 的上同调维数为 2，且满足 Euler–Poincaré 维数的消没。这意味着
\[
\dim_L T_xX_L=\dim H_L+\dim_L H^2(W^0_t,\mathrm{Lie}\,H_L).
\]
为计算 $H^2(W^0_t,\mathrm{Lie}\,H_L)=H^1(\langle\mathrm{Fr}\rangle,H^1(s^{\Z[\frac1q]},\mathrm{Lie}\,H_L))$，我们首先注意典范映射 $H^1(s^{\Z},\mathrm{Lie}\,H_L)=(\mathrm{Lie}\,H_L)_\sigma\to H^1(s^{\Z[\frac1q]},\mathrm{Lie}\,H_L)$ 是同构（这里 $\sigma=\varphi(s)$，下标 $\sigma$ 表示 $\sigma$-余不变量）。确实，由于 $\sigma_{\mathrm{ss}}$ 的阶有限，所有 $\varphi(s^{q^{-n}})$ 都属于 $\mathrm{Aut}_L(\mathrm{Lie}\,H_L)$ 中 $\langle\sigma\rangle$ 的 Zariski 闭包。经由这个同构，$\mathrm{Fr}$ 在 $H^1(\sigma^{\Z[\frac1q]},\mathrm{Lie}\,H_L)$ 上的作用\tn{底本此处作「$H^1(\sigma^{\Z[1/q]},\mathrm{Lie}\,H_L)$」，上文同一上同调群写作 $s^{\Z[1/q]}$，此处 $\sigma$ 疑为 $s$ 之误。}对应于 $F=\varphi(\mathrm{Fr})$ 在 $(\mathrm{Lie}\,H_L)_\sigma$ 上的自然作用乘以 $q^{-1}$。于是得到同构 $H^2(W^0_t,\mathrm{Lie}\,H_L)\simeq(\omega^{-1}\otimes(\mathrm{Lie}\,H_L)_\sigma)_F$，而在取对偶之后我们看到
\[
H^2(W^0_t,\mathrm{Lie}\,H_L)^*\simeq H^0\bigl(W^0_t,(\mathrm{Lie}\,H_L)^*\otimes\omega\bigr).
\]
\end{proof}

\ysrc{14}现在我们要证明 $X_L$ 一般光滑，至少在 $L$ 的特征为 $0$ 或足够大时如此。下一个引理处理特征 $0$ 的情形。关于结果在正特征下仍然成立的讨论，参见 [DHKM20, §5]。

\begin{prop}{3.7}
若 $\mathrm{Char}\,L=0$，则 $X_L$ 的任一不可约分支都含有光滑点。因此 $X_{H,1}$ 一般光滑，特别地是约化的。
\end{prop}

\begin{proof}
由注 3.4，只需证明：对任意 $(F,\sigma)\in X(L)$，存在 $c\in(H_\sigma)^0(L)$ 使得 $(Fc,\sigma)$ 是光滑点。由上一个引理，只需找 $c\in(H_\sigma)^0(L)$ 使 $Fc$ 在 $\bigl((\mathrm{Lie}\,H_L)^*\otimes\omega\bigr)^\sigma$ 中没有不动向量，也就是说，使 $q^{-1}$ 不是 $Fc$ 在 $(\mathrm{Lie}\,H_L)^{*,\sigma}$ 上的特征值。

先假设 $\sigma$ 幂幺（从而属于 $H^0(L)$）。由 Jacobson–Morozov 定理，可找到余特征标 $\lambda:\mathbb G_m\to H$，使得 $\log\sigma$ 在 $\mathrm{Ad}_\lambda$ 作用于 $\mathrm{Lie}\,H$ 时具有权 2。此外这样的 $\lambda$ 在幂幺根 $R_u(H_\sigma)$ 的元素的共轭下唯一。于是存在 $c_1\in R_u(H_\sigma)(L)$ 使 $Fc_1$ 中心化 $\lambda$。现在置 $F_\lambda:=\lambda(q^{\frac12})\in H(L)$。有 $F_\lambda\sigma F_\lambda^{-1}=\sigma^q$，故 $F_\lambda^{-1}Fc_1\in H_\sigma\cap H_\lambda$。已知代数群 $H_\sigma\cap H_\lambda$ 是 $H_\sigma$ 的 Levi 补（见如 [Jan04, Prop. 5.10]）。特别地它是约化的，故可固定其单位分支 $(H_\sigma\cap H_\lambda)^0$ 的一个固定数据，并找到 $c_2\in(H_\sigma\cap H_\lambda)^0$ 使 $F_\lambda^{-1}Fc_1c_2$ 正规化该固定数据。置 $m:=|\pi_0(H_\sigma\cap H_\lambda)|$，这表明 $(F_\lambda^{-1}Fc_1c_2)^m$ 属于中心 $Z(H_\sigma\cap H_\lambda)$，甚至属于在 $F_\lambda^{-1}Fc_1c_2$ 的共轭下不变的子群 $Z(H_\sigma\cap H_\lambda)^{F_\lambda^{-1}Fc_1c_2}$。若 $m'$ 记 $\pi_0\bigl(Z(H_\sigma\cap H_\lambda)^{F_\lambda^{-1}Fc_1c_2}\bigr)$ 的阶，我们看到 $(F_\lambda^{-1}Fc_1c_2)^{mm'}$ 是环面 $\bigl(Z(H_\sigma\cap H_\lambda)^{F_\lambda^{-1}Fc_1c_2}\bigr)^0$ 的一个元素。由于环面是可除群，存在 $c_3\in\bigl(Z(H_\sigma\cap H_\lambda)^{F_\lambda^{-1}Fc_1c_2}\bigr)^0$ 使得 $(F_\lambda^{-1}Fc_1c_2c_3)^{mm'}=1$。置 $c:=c_1c_2c_3\in(H_\sigma)^0$ 与 $c':=F_\lambda^{-1}Fc\in H_\sigma\cap H_\lambda$，我们看到 $Fc=F_\lambda c'=c'F_\lambda$，从而 $Fc$ 在 $(\mathrm{Lie}\,H_L)^{*,\sigma}$ 上的特征值是 $c'$ 与 $F_\lambda$ 的特征值之积。$c'$ 的特征值是单位根。另一方面，有分解 $(\mathrm{Lie}\,H_L)^*=\mathrm{Lie}(Z(H)_L)^*\oplus(\mathrm{Lie}\,H_{\mathrm{der},L})^*$。$F_\lambda$ 在 $\mathrm{Lie}(Z(H)_L)^*$ 上的特征值仍是单位根，并且我们知道 $\lambda$ 在 $\mathrm{Lie}\,H_\sigma$ 上具有非负的权。由于 Killing 型把 $(\mathrm{Lie}\,H_{\mathrm{der},L})^*$ 与 $\mathrm{Lie}\,H_{\mathrm{der},L}$ 等同，我们看到 $q^{-1}$ 不是 $Fc$ 在 $(\mathrm{Lie}\,H_L)^{*,\sigma}=\mathrm{Lie}\,H_\sigma$ 上的特征值。

现在把一般情形归约到幂幺情形。设 $\sigma=\sigma_{\mathrm{ss}}\sigma_{\mathrm{u}}$ 是 $\sigma$ 的 Jordan 分解。由于 $F$ 正规化约化代数子群 $(H_{\sigma_{\mathrm{ss}}})^0$，我们可以重复上面的论证，找到 $h\in(H_{\sigma_{\mathrm{ss}}})^0$ 使 $Fh$ 正规化该群的一个固定数据且阶有限。这意味着由 $H_{\sigma_{\mathrm{ss}}}$ 与 $F$ 生成的子群 $H'$ 的 Zariski 闭包具有单位分支 $(H_{\sigma_{\mathrm{ss}}})^0$。由幂幺情形，存在 $c\in(H'_{\sigma_{\mathrm{u}}})^0=(H_\sigma)^0$ 使 $q^{-1}$ 不是 $Fc$ 在 $\mathrm{Lie}\,H'_{\sigma_{\mathrm{u}}}=\mathrm{Lie}\,H_\sigma$ 上的特征值，此即所求。
\end{proof}

\subsection{忠实平坦性与 $\ell$-进连续性}

下一个陈述蕴含：$X_{H,1}$ 的所有不可约分支（带其约化子概形结构）在 $\Lambda$ 上忠实平坦；并且 $X_{H,1}$ 的函数环 $R_{H,1}$ 对所有素数 $\ell\ne p$ 都是 $\ell$-进分离的。

\begin{lem}{3.8}
$X_{H,1}$ 的任一不可约分支都含有某个 $\overline\Z_\ell$-点，对任意素数 $\ell\ne p$。特别地，$X_{H,1}$ 的函数环 $R_{H,1}$ 是 $\ell$-进分离的。
\end{lem}

\begin{proof}
设 $Y$ 是 $X=X_{H,1}$ 的一个不可约分支，并选取一个嵌入 $\Lambda\hookrightarrow\overline\Z_\ell$。由于 $X_{H,1}$ 在 $\Lambda$ 上平坦，基变换 $Y_{\overline\Q_\ell}$ 非空，
